\section{Constructive complex numbers}
\label{sec:complex}

\subsection{Carrier and operations}
Let $R=\mathbb{R}_{\mathrm H}$.  The formal carrier is the product
$\CH=R\times R$.  Its operations are
\begin{align}
0_{\CH}&=(0,0), & 1_{\CH}&=(1,0),\\
(a,b)+(c,d)&=(a+c,b+d), & -(a,b)&=(-a,-b),\\
(a,b)(c,d)&=(ac-bd,ad+bc),\\
i&=(0,1), & \overline{(a,b)}&=(a,-b).
\end{align}
Thus the imaginary unit and conjugation are checked definitions, although no
separate theorem about $i^2$ or conjugation multiplicativity is claimed by
these modules.

\formal{\texttt{src/ConstructiveComplex/Complex.agda}}
{\texttt{Complex}; \texttt{complexZero}; \texttt{complexOne}; \texttt{complexAdd}; \texttt{complexNeg}; \texttt{complexMul}; \texttt{imaginaryUnit}; \texttt{complexConjugate}}
{These are definitions, with the displayed coordinate formula for multiplication.}

The only HoTT Reals declarations used directly are the carrier
\agda{HoTTReals.Data.Real.Base.ℝ}, the set proof
\agda{RealProperties.ℝ-isSet}, and the commutative-ring package
\agda{RealMultiplication.ℝCommutativeRing}.  The local development does not
reproduce Brough's construction of real numbers.

\upstream{\texttt{HoTTReals/Data/Real/Base.agda}}{\texttt{ℝ}}
\upstream{\texttt{HoTTReals/Data/Real/Properties.agda}}{\texttt{ℝ-isSet}}
\upstream{\texttt{HoTTReals/Data/Real/Algebra/Multiplication.agda}}{\texttt{ℝCommutativeRing}}

\subsection{Sethood and pair paths}
Since $R$ is a set, its product is a set by the imported closure theorem
\agda{isSet×}.  Equalities of complex numbers are assembled from equalities
of real coordinates with \agda{ΣPathP}:
\begin{equation}
 a=b\ \text{and}\ c=d\quad\Longrightarrow\quad(a,c)=(b,d).
\end{equation}
\formal{\texttt{src/ConstructiveComplex/Complex.agda}}
{\texttt{complexIsSet}; \texttt{pairPath}}
{\texttt{complexIsSet} uses \texttt{isSet×}; the private helper \texttt{pairPath} is \texttt{ΣPathP}.}

\subsection{The ring proof}
\begin{proposition}
The displayed operations make $\CH$ a commutative ring.
\end{proposition}
\begin{proof}
All eight algebraic laws reduce to polynomial identities in the commutative
ring $R$.  For example, associativity compares
\begin{align*}
(a,b)((c,d)(e,f))&=\bigl(a(ce-df)-b(cf+de),\ a(cf+de)+b(ce-df)\bigr),\\
((a,b)(c,d))(e,f)&=\bigl((ac-bd)e-(ad+bc)f,\ (ac-bd)f+(ad+bc)e\bigr).
\end{align*}
The two coordinates agree by associativity, commutativity, distributivity,
and additive-inverse laws in $R$; \agda{ΣPathP} then gives equality in
$\CH$.  Addition, the unit and inverse laws, distributivity, and
commutativity of multiplication are proved by the same coordinatewise
method.  The sethood proof and these eight fields are precisely the data
accepted by Cubical's commutative-ring constructor.
\end{proof}

The formal proof discharges each coordinate by the verified commutative-ring
solver instantiated at Brough's ring, rather than treating the displayed
calculation as an unchecked external argument:
\begin{verbatim}
complexMulAssociative (a , b) (c , d) (e , f) =
  pairPath (solve! RealMultiplication.ℝCommutativeRing)
           (solve! RealMultiplication.ℝCommutativeRing)
\end{verbatim}
\formal{\texttt{src/ConstructiveComplex/Complex.agda}}
{\texttt{complexAddAssociative}; \texttt{complexAddUnitRight}; \texttt{complexAddInverseRight}; \texttt{complexAddCommutative}; \texttt{complexMulAssociative}; \texttt{complexMulUnitRight}; \texttt{complexMulDistributesOverAddLeft}; \texttt{complexMulCommutative}; \texttt{complexIsCommRing}; \texttt{ComplexCommRing}}
{The eight local laws are combined by \texttt{makeIsCommRing}; \texttt{ComplexCommRing} is the resulting level-zero package.}
